\documentclass[10pt]{article}
\usepackage[T1]{fontenc}
\usepackage[utf8]{inputenc}
\usepackage[french]{babel}
\usepackage{amsmath,amssymb,mathrsfs,mathpazo}
\usepackage{geometry}
\geometry{a4paper,margin=2cm}
\usepackage{graphicx,array,tabularx,tabularray,colortbl,multirow,multicol,booktabs,enumitem,pifont,ulem,fancybox,diagbox,cancel}
\usepackage{pstricks,pstricks-add,pst-plot,pst-tree,pst-node,pst-text,pst-func,pst-math,pst-bspline,pst-3dplot,pst-eucl,pst-solides3d}
\let\pscalculate\undefined
\usepackage{pst-bezier}
\usepackage{tikz,pgfplots,tkz-tab}
\usetikzlibrary{arrows,arrows.meta,patterns,calc,positioning,trees,decorations.markings,shapes.geometric,shapes.misc,shapes.arrows,intersections,decorations.pathmorphing,decorations.pathreplacing,backgrounds,shadows,babel}
\pgfplotsset{compat=1.18}
\usepgfplotslibrary{fillbetween}
\usepackage[np]{numprint}
\usepackage{listings,makecell,pas-tableur,twoopt}
\let\barre\undefined
\usepackage[pstricks]{tablvar}
\definecolor{darkgreen}{RGB}{0,100,0}
\providecommand{\red}{\color{red}}
\providecommand{\blue}{\color{blue}}
\providecommand{\green}{\color{green}}
\providecommand{\black}{\color{black}}
\providecommand{\white}{\color{white}}
\providecommand{\gray}{\color{gray}}
\providecommand{\cyan}{\color{cyan}}
\providecommand{\magenta}{\color{magenta}}
\providecommand{\yellow}{\color{yellow}}
\providecommand{\violet}{\color{violet}}
\providecommand{\orange}{\color{orange}}
\providecommand{\euro}{\texteuro}
\providecommand{\EUR}{\texteuro}
\providecommand{\ieme}{\textsuperscript{e}}
\providecommand{\degres}{\ensuremath{^\circ}}
\providecommand{\degree}{\ensuremath{^\circ}}
\providecommand{\vv}[1]{\overrightarrow{#1}}
\providecommand{\Coord}[1]{#1}
\providecommand{\decosix}{}
\providecommand{\starredbullet}{\ensuremath{\bullet}}
\pagestyle{empty}
\setlength{\parindent}{0pt}
\usepackage{adjustbox}\DeclareUnicodeCharacter{2236}{:}\DeclareUnicodeCharacter{2190}{\ensuremath{\leftarrow}}
\providecommand{\R}{\mathbb{R}}
\providecommand{\N}{\mathbb{N}}
\providecommand{\Z}{\mathbb{Z}}
\providecommand{\Q}{\mathbb{Q}}
\providecommand{\D}{\mathbb{D}}
\providecommand{\C}{\mathbb{C}}
\providecommand{\e}{{\mathrm{e}}}
\providecommand{\vect}[1]{\overrightarrow{#1}}
\providecommand{\vectt}[1]{\overrightarrow{\mathrm{#1}}}
\providecommand{\vvt}[1]{\overrightarrow{\mathrm{#1}}}
\providecommand{\np}[2][]{{#2~#1}}
\providecommand{\vv}[1]{{\overrightarrow{#1}}}
\renewcommand{\d}{{\mathrm{d}}}
\providecommand{\degres}{{^\circ}}
\providecommand{\degre}{{^\circ}}
\providecommand{\blue}{\color{blue}}
\providecommand{\red}{\color{red}}
\providecommand{\enskip}{\quad}
\providecommand{\llbracket}{\lbrack\!\lbrack}
\providecommand{\rrbracket}{\rbrack\!\rbrack}
\providecommand{\euro}{{\text{€}}}
\providecommand{\no}{{\text{n°}}}
\providecommand{\interval}[2]{{[#1;#2]}}
\providecommand{\intervFO}[2]{{[#1;#2[}}
\providecommand{\dotfill}{\ldots}

\providecommand{\ds}{\displaystyle}
\providecommand{\textding}[1]{#1}
\providecommand{\Oij}{\(\mathrm{O};\vec{\imath},\vec{\jmath}\)}
\providecommand{\Oijk}{\(\mathrm{O};\vec{\imath},\vec{\jmath},\vec{k}\)}
\providecommand{\Ouv}{\(\mathrm{O};\vec{u},\vec{v}\)}
\providecommand{\intervFF}[2]{\ensuremath{[#1;#2]}}
\renewcommand{\euro}{\ifmmode\text{\texteuro}\else\texteuro\fi}
\DeclareUnicodeCharacter{00B0}{\ifmmode^{\circ}\else\textdegree\fi}
\DeclareUnicodeCharacter{20AC}{\ifmmode\text{\texteuro}\else\texteuro\fi}
% Énoncés : Éducation nationale. Transcription et corrigés : APMEP, Denis Vergès et contributeurs. Découpage et classement : CoopMaths, Sébastien Lozano, Éric Elter et contributeurs. Adaptation de la présentation : Maths Explorer.
% CC BY-SA : https://coopmaths.fr/www/about/
% https://coopmaths.fr/alea/static/bac/2022/tex/bac_2022_05_sujet2_ameriquenord_1_cor.tex
\begin{document}
\noindent\textit{APMEP / CoopMaths. CC BY-SA. Présentation adaptée par Maths Explorer.}\par\medskip
\medskip

\begin{enumerate}
	\item L'arbre complété avec les valeurs disponibles :
		\begin{center}
		\pstree[treemode=R,nodesepA=0pt,nodesepB=2.5pt,treesep=1cm,levelsep=2.5cm]{\TR{}}
		{\pstree{\TR{$A_1$~~}\taput{0,5}}
			{\TR{$A_2$}\taput{0,84}
			\TR{$B_2$}\tbput{0,16}
			}
		\pstree{\TR{$B_1$~~}\tbput{0,5}}
			{\TR{$A_2$}\taput{0,24}
			\TR{$B_2$}\tbput{0,76}
			}
		}
		\end{center}
	\item 
	\begin{enumerate}
		\item 	Utilisons la formule des probabilités totales pour calculer $a_2=p(A_2)$ :
		
		$\begin{gathered}a_2=p(A_2\cap A_1)+p(A_2\cap B_1)=p_{A_1}(A_2)\times p(A_1)+p_{B_1}(A_2)\times p(B_1)\\=0,84\times 0,5+0,24\times 0,5=0,54\end{gathered}$. Donc $a_2=0,54$.
		\item Utilisons la formule de Bayes pour calculer $p_{A_2}(B_1)$ :
		
		$p_{A_2}(B_1)=\dfrac{p(A_2\cap B_1)}{p(A_2)}=\dfrac{p_{B_1}(A_2)\times p(B_1)}{p(A_2)}=\dfrac{0,24\times 0,5}{0,54}\approx 0,222$ 
	\end{enumerate}
	\item 
	\begin{enumerate}
		\item On remarquera au préalable que, $\forall n\in \N^*$, $a_n+b_n=1$.
		
		L'arbre complété avec les valeurs disponibles :
			\begin{center}
				\pstree[treemode=R,nodesepA=0pt,nodesepB=2.5pt,treesep=1cm,levelsep=2.5cm]{\TR{}}
				{\pstree{\TR{$A_n$~~}\taput{$a_n$}}
					{\TR{$A_{n+1}$}\taput{0,84}
					\TR{$B_{n+1}$}\tbput{0,16}
					}
				\pstree{\TR{$B_n$~~}\tbput{$1-a_n$}}
					{\TR{$A_{n+1}$}\taput{0,24}
					\TR{$B_{n+1}$}\tbput{0,76}
					}
				}
			\end{center} 
		\item Utilisons là encore, la formule des probabilités totales pour déterminer $a_{n+1}$ en fonction de $a_n$, pour tout entier naturel non nul :
		
		$\forall n\in \N^*$, $\begin{gathered}a_{n+1}=p(A_{n+1}\cap A_n)+p(A_{n+1}\cap B_n)=p_{A_n}(A_{n+1})\times p(A_n)+p_{B_n}(A_{n+1})\times p(B_n)\\=0,84\times p(A_n)+0,24\times p(B_n)=0,84~a_n+0,24~b_n\end{gathered}$. Or $\forall n\in \N^*$, $b_n=1-a_n$.
		
		Donc $\forall n\in \N^*$, $a_{n+1}=0,84~a_n+0,24(1-a_n)=0,6~a_n+0,24$
		\item Montrons par récurrence que $\forall n \in \N^*$, $a_n = 0,6 - 0,1 \times 0,6^{n - 1}$.
						
		\textit{Initialisation :} $a_1=0,6 - 0,1 \times 0,6^{1 - 1}=0,6-0,1=0,5$. L'initialisation est vérifiée.
				
		\textit{Hérédité :} Soit  $n \in \N^*$, et supposons que $a_n = 0,6 - 0,1 \times 0,6^{n - 1}$.
		
		Montrons que $a_{n+1} = 0,6 - 0,1 \times 0,6^{n}$.
				
		D'après la question précédente,$a_{n+1}=0,6~a_n+0,24$, donc en utilisant l'hypothèse de récurrence,
		
		$a_{n+1}=0,6\left(0,6 - 0,1 \times 0,6^{n - 1}\right)+0,24=0,36-0,1\times0,6 \times0,6^{n - 1}+0,24=0,6-0,1\times 0,6^n $. On obtient ce qu'il fallait démontrer. L'hérédité est démontrée.
				
		\textit{Conclusion : }  La proposition est vraie au rang $0$ et si elle est vraie au rang $n$, elle l'est aussi au rang $n + 1$.	D'après l'axiome de récurrence, $\forall n \in \N^*$, $a_n = 0,6 - 0,1 \times 0,6^{n - 1}$.
		\item $\displaystyle\lim\limits_{n \rightarrow +\infty} 0,6^{n-1}=0$ car $0,6\in ]-1~;~1[$. Donc $\displaystyle\lim\limits_{n \rightarrow +\infty} a_n=0,6$.
		
		Cela signifie qu'au bout d'un certain temps, la probabilité qu'un vélo soit à la station A est de 60 \%.
		\item Résolvons : $a_n\geqslant 0,599$.
		$\begin{gathered}a_n\geqslant 0,599\iff  0,6 - 0,1 \times 0,6^{n - 1} \geqslant 0,599 \\
		\iff - 0,1 \times 0,6^{n - 1} \geqslant -0,001 \iff 0,6^{n - 1} \leqslant \dfrac{-0,001}{-0,1} \iff 0,6^{n - 1} \leqslant \dfrac{1}{100} \end{gathered}$.
		
		Sachant que la fonction $x\mapsto\ln(x)$ est strictement croissante sur $\R_+^*$, 
		
		$0,6^{n - 1} \leqslant \dfrac{1}{100} \iff \ln\left(0,6^{n - 1} \right) \leqslant \ln\left(\dfrac{1}{100}\right) \iff (n-1)\times \ln(0,6) \leqslant -\ln(100)$. Or $\ln(0,6)<0$, donc
		$n-1 \geqslant \dfrac{-\ln(100)}{\ln(0,6)} \iff n \geqslant 1-\dfrac{2\ln(10)}{\ln(0,6)}$.
		
		À la calculatrice, $1-\dfrac{2\ln(10)}{\ln(0,6)}\approx 10,02$ donc $n\geqslant 11$.
		
		La probabilité que le vélo se trouve au point A est supérieure à $0,599$ à partir du 11-ième jour.
		\end{enumerate}
\end{enumerate}

\bigskip

\end{document}
